% part: intuitionistic-logic
% chapter: propositions-as-types
% section: normalization

\documentclass[../../../include/open-logic-section]{subfiles}

\begin{document}

\olfileid{int}{pty}{nor}

\olsection{Normalisasi}

Ing bagean iki kita mbuktekake yen, kanthi sawijining urutan reduksi,
saben dedhuksi bisa direduksi dadi dedhuksi normal, sing diarani
\emph{sipat normalisasi}. Kita bakal nggunakake korespondhensi
proposisi minangka tipe: kita nuduhake yen saben terma bukti sing tipene bener bisa
direduksi dadi wujud normal; banjur normalisasi tumrap
!!{derivation}s dedhuksi alami bisa dipikolehi.

Sepisanan kita netepake sawetara fungsi kanggo ngukur kerumitan terma.
\emph{Dawa}~$\len{!A}$ sawijining !!a{formula}s ditetepake kanthi
\begin{align*}
  \len{p} &= 0\\
  \len{\lfalse} &= 0\\
  \len{!A \land !B} &= \len{!A} + \len{!B} + 1\\
  \len{!A \lor !B} &= \len{!A} + \len{!B} + 1 \\
  \len{!A \lif !B} &= \len{!A} + \len{!B} + 1.
\end{align*}
Kerumitan redex~$M$ diukur kanthi \emph{rangking
  cut}~$\cutrank{M}$:
\begin{align*}
  \cutrank{(\lambd[\typeof{x}{!A}][\typeof{N}{!B}])Q} &=
  \len{!A} + \len{!B} + 1 \\
  \cutrank{\ande{i}{\andi{\typeof{M}{!A}}{\typeof{N}{!B}}}} &=
  \len{!A} + \len{!B} + 1 \\
  \cutrank{\ore{\ori{i}{!A_{i}}{\typeof{M}{!A_i}}}
    {\typeof{x_1}{!A_1}}{\typeof{N_1}{!C}}
    {\typeof{x_2}{!A_2}}{\typeof{N_2}{!C}}} &=
  \len{!A_1} + \len{!A_2} + 1
\end{align*}
Kerumitan terma bukti diukur kanthi redex sing paling ruwet ing terma
kasebut, lan regane $0$ yen terma iku normal:
\[
  \maxrank{M} = \max (\{0\} \cup \{\cutrank{N} | N \text{ iku subterma saka } M \text{
    lan minangka redex}\})
\]

\begin{lem}\ollabel{lem:subst}
  Yen $\Subst{M}{\typeof{N}{!A}}{\typeof{x}{!A}}$ iku redex lan $M
  \not\ident x$, salah siji kasus ing ngisor iki dumadi:
  \begin{enumerate}
  \item $M$ dhewe iku redex, utawa
  \item $M$ awujud $\ande{i}{x}$, lan $N$ awujud
    $\andi{P_1}{P_2}$
  \item $M$ awujud $\ore{x}{x_1}{P_1}{x_2}{P_2}$, lan $N$
    awujud $\ori{i}{}{Q}$
  \item $M$ awujud $x Q$, lan $N$ awujud
    $\lambd[x][P]$
  \end{enumerate}
  Ing kasus kapisan, $\cutrank{\Subst{M}{N}{x}} = \cutrank{M}$; ing
  kasus liyane, $\cutrank{\Subst{M}{N}{x}} = \len{!A}$.
\end{lem}
\begin{proof}
  Bukti kanthi induksi marang~$M$.
  \begin{enumerate}
  \item Yen $M$ mung variabel $y$ lan $y \not\ident x$,
    $\Subst{y}{N}{x}$ iku $y$, mula dudu redex.
  \item Yen $M$ awujud $\andi{N_1}{N_2}$, utawa $\lambd[x][N]$, utawa
    $\ori{i}{!A}{N}$, $\Subst{M}{\typeof{N}{!A}}{\typeof{x}{!A}}$ uga
    nduweni wujud kasebut, mula dudu redex.
  \item Yen $M$ awujud $\ande{i}{P}$, kita nimbang rong kasus.
    \begin{enumerate}
      \item Yen $P$ awujud $\andi{P_1}{P_2}$, $M \ident
        \ande{i}{\andi{P_1}{P_2}}$ iku redex, lan cetha yen
        \[
        \Subst{M}{N}{x} \ident
        \ande{i}{\andi{\Subst{P_1}{N}{x}}{\Subst{P_2}{N}{x}}}
        \]
        uga redex. Rangking cut kalorone padha.
      \item Yen $P$ mung variabel, variabel iku kudu $x$ supaya asil
        substitusi dadi redex, lan $N$ kudu awujud
        $\andi{P_1}{P_2}$. Saiki gatekna
        \[
        \Subst{M}{N}{x} \ident \Subst{\ande{i}{x}}{\andi{P_1}{P_2}}{x},
        \]
        sing padha karo $\ande{i}{\andi{P_1}{P_2}}$. Rangking cute padha karo
        dawa tipe variabel sing diganti, yaiku $\len{!A}$.
    \end{enumerate}
  \end{enumerate}
  Kasus $\ore{N}{x_1}{N_1}{x_2}{N_2}$ lan $P Q$ padha carane.
\end{proof}

\begin{lem}
  Yen $M$ dikontraksi dadi $M'$, lan $\cutrank{M} > \cutrank{N}$ kanggo
  saben subterma redex sejati $N$ saka $M$, mula $\cutrank{M} >
  \maxrank{M'}$.
\end{lem}

\begin{proof}
  Bukti kanthi pamisahan kasus.
  \begin{enumerate}
  \item Yen $M$ awujud $\ande{i}{\andi{M_1}{M_2}}$, $M'$ iku
    $M_i$; amarga saben subterma saka $M_i$ uga subterma sejati
    saka $M$, pratelan kasebut bener.
  \item Yen $M$ awujud
    $(\lambd[\typeof{x}{!A}][N])\typeof{Q}{!A}$, $M'$ iku
    $\Subst{N}{\typeof{Q}{!A}}{\typeof{x}{!A}}$. Gatekna redex ing
    $M'$. Bisa uga ana redex sing cocog ing $N$ kanthi rangking cut
    sing padha, sing luwih cilik tinimbang $\cutrank{M}$ miturut asumsi,
    utawa redex iku disalin saka argumen $Q$ lan rangking cute luwih cilik
    miturut asumsi sing padha, utawa rangking cute padha karo $\len{!A}$,
    sing miturut dhefinisi luwih cilik tinimbang
    $\cutrank{(\lambd[x^{!A}][N])Q}$.
  \item Yen $M$ awujud
    \[
    \ore{\ori{i}{}{\typeof{N}{!A_i}}}
        {\typeof{x_1}{!A_1}}{\typeof{N_1}{!C}}
        {\typeof{x_2}{!A_2}}{\typeof{N_2}{!C}},
    \]
    mula $M' \ident \Subst{N_i}{N}{\typeof{x_i}{!A_i}}$. Gatekna
    redex ing~$M'$. Bisa uga ana redex sing cocog ing $N_i$ kanthi
    rangking cut sing padha, sing luwih cilik tinimbang $\cutrank{M}$ miturut
    asumsi; utawa redex iku disalin saka terma injeksi $N$ lan rangking cute
    luwih cilik miturut asumsi sing padha; utawa rangking cute padha karo
    $\len{!A_i}$, sing miturut dhefinisi luwih cilik tinimbang
    $\cutrank{\ore{\ori{i}{}{\typeof{N}{!A_i}}}
      {\typeof{x_1}{!A_1}}{\typeof{N_1}{!C}}
      {\typeof{x_2}{!A_2}}{\typeof{N_2}{!C}}}$.
  \end{enumerate}
\end{proof}

\begin{thm}
  Saben terma bukti sing tipene bener direduksi dadi wujud normal; saben !!{derivation}s
  direduksi dadi !!{derivation}s normal.
\end{thm}

\begin{proof}
  Pratelan kapindho minangka akibat saka pratelan kapisan. Kita mbuktekake
  pratelan kapisan kanthi induksi lengkap marang $m=\maxrank{M}$,
  kanthi $M$ minangka terma bukti sing tipene bener.
  \begin{enumerate}

  \item Yen $m=0$, $M$ wis normal.
  \item
    Yen ora, kita nggunakake induksi marang~$n$, yaiku cacah redex
    ing $M$ sing rangking cute padha karo~$m$.
    \begin{enumerate}
    \item Yen $n=1$, pilih redex $N$ kanthi $m = \cutrank{N} >
      \cutrank{P}$ kanggo saben subterma sejati $P$ sing uga redex.
      Redex kaya mangkono mesthi ana, amarga saben terma mung nduweni
      subterma sing cacahe winates.

      Upamane $N'$ nuduhake asil kontraksi $N$. Miturut lema,
      $\maxrank{N'} < \maxrank{N}$, mula kita bisa ndeleng yen $n$,
      yaiku cacah redex kanthi $\cutrank=m$, suda. Mula $m$ suda
      (paling ora $1$), lan kita bisa nggunakake hipotesis induksi
      kanggo~$m$.
    \item Kanggo langkah induksi, upamane $n>1$. Prosese padha,
      nanging $n$ mung suda dadi cacah positif, mula $m$ ora owah.
      Kita banjur nggunakake hipotesis induksi kanggo~$n$.
    \end{enumerate}
\end{enumerate}
\footnote{Cathetan editor (OLPL-655): Argumen sumber iki mung matesi redex
ing subterma sing dikontraksi. Argumen iki durung ngendhaleni redex anyar
sing kawiyak ing konteks saubenge, mligine nalika eliminasi disjungsi
nduweni tipe asil sembarang. Rangking sing ditetepake uga durung menehi
ukuran tumrap konversi permutasi. Mula argumen sing dicithak iki isih
durung jangkep kanggo kalkulus sakabehe; cathetan iki ora menehi bukti
sing isih kurang kasebut.}
\end{proof}

Kasunyatan yen terma $\lambd$ mawa tipe bisa dinormalisasi tegese saben
terma kasebut \emph{nduweni wujud normal}. Yen kita nganggep reduksi
$\beta$ tumrap terma $\lambd$ menyang wujud normale minangka nindakake
komputasi, lan wujud normal minangka nilai komputasi kasebut, iki tegese
saben komputasi ing kalkulus $\lambd$ mawa tipe pungkasane mandheg lan
ngasilake nilai. Kajaba iku, panglestaren tipe njamin yen nilai kasebut
nduweni tipe sing bener. Contone, yen $N$ nduweni tipe $!A \lif !B$
(yaiku fungsi sing dikarepake bisa nampa argumen mawa tipe~$!A$ lan
ngasilake nilai mawa tipe~$!B$), lan kita ngetrapake fungsi iku marang
terma~$M$ mawa tipe~$!A$ (input), $(NM)$ direduksi dadi terma $N'$
(output), lan output iku mesthi mawa tipe~$!B$.

Sejatine, terma ing kalkulus $\lambd$ mawa tipe ora mung direduksi dadi
wujud normal kanthi siji cara ngetrapake reduksi $\beta$ marang
subterma, nanging \emph{saben} cara ngetrapake reduksi $\beta$ marang
subterma pungkasane mandheg ing wujud normal. Sipat iki diarani
\emph{normalisasi kuwat}. Tegese, ora mung ana cara kanggo njamin yen
komputasi pungkasane mandheg lan ngasilake nilai, nanging \emph{saben}
cara nindakake komputasi pungkasane ngasilake nilai.

Kepiye kita ngerti yen cara-cara nindakake komputasi sing beda ngasilake
nilai sing \emph{padha}? Nganti saiki kita mung ngerti (saka panglestaren
tipe) yen kabeh nilai sing diasilake nduweni tipe sing padha. Kalkulus
$\lambd$ mawa tipe nduweni sipat wigati liyane: kalkulus iki
\emph{konfluen}, utawa \emph{Church-Rosser}. Yen $N \red N_1$ lan $N \red
N_2$, ana terma~$N'$ kanthi $N_1 \red N'$ lan $N_2 \red N'$. Elinga
yen $N \red N'$ tegese $N'$ diasilake saka nol utawa luwih langkah
reduksi $\redone$. Yen $N_1$ awujud normal, siji-sijine terma $N'$
kanthi $N_1 \red N'$ iku $N_1$ dhewe (kanthi nol langkah reduksi
$\redone$). Mula yen $N_1$ lan $N_2$ kalorone awujud normal, kita
nduweni $N_1 = N' = N_2$. Kanthi tembung liya, amarga $\red$
nduweni normalisasi kuwat, saben cara ngreduksi terma pungkasane
ngasilake wujud normal. Amarga $\red$ nduweni sipat Church-Rosser,
saben rong wujud normal padha. Mula komputasi ing kalkulus $\lambd$
mawa tipe tansah mandheg, lan nilai (wujud normal) sing diasilake
padha, sanajan cara komputasine beda.

Mbuktekake yen sawijining relasi konfluen lumrahe angel. Nanging,
kita bisa netepake syarat sing luwih ringkih ing ngisor iki:

\begin{prop}[Sipat Church-Rosser Ringkih]
Yen $N \redone N_1$ lan $N \redone N_2$, ana terma $N'$ kanthi
$N_1 \red N'$ lan $N_2 \red N'$.
\end{prop}

\begin{proof}
  Rong terma $N_1$ lan $N_2$ diasilake saka $N$ kanthi ngganti redex
  $M_1$ utawa $M_2$ nganggo asil kontraksine. Redex $M_1$ lan $M_2$
  iku subterma saka $N$, mula ora bisa tumpang tindih mung saweneh.
  Tegese, salah siji iku subterma saka sijine, utawa kalorone dumunung
  ing papan beda sing ora tumpang tindih ing $N$. Gatekna kasus kapindho:
  $N$ awujud $N(M_1,M_2)$. Ngganti $M_1$ nganggo asil kontraksine $M_1'$
  ngasilake $N_1$, yaiku $N(M_1',M_2)$, lan ngganti $M_2$ nganggo
  asil kontraksine $M_2'$ ngasilake $N_2$, yaiku $N(M_1,M_2')$.
  Kalorone direduksi dadi $N(M_1',M_2')$. Kanggo kasus liyane,
  upamane $M_2$ iku subterma saka~$M_1$ (kasus kosokbaline simetris),
  lan bedakna kasus miturut jinis redex $M_1$. Contone, yen $M_1 \ident
  \proj{1}{\pair{O_1}{O_2}}$, terma iki direduksi dadi $O_1$. Yen $M_2$
  iku subterma saka $O_1$, ngowahi $M_2$ ngasilake $O_1'$. Mula
  ngowahi $M_1$ dhisik, banjur $M_2$, ngasilake $O_1'$. Nanging,
  ngowahi $M_2$ dhisik ngasilake $\proj{1}{\pair{O_1'}{O_2}}$, lan
  iki uga direduksi dadi~$O_1'$. Yen redex njero dumunung ing komponen
  sing dibuwang, loro urutan kasebut tetep ngasilake komponen kapisan.
  \footnote{Cathetan editor (OLPL-657): Iki minangka sketsa kasus kontraksi
  utama. Kasus njero liyane lan sesambungane karo konversi permutasi
  durung dibuktekake ing kene.}
\end{proof}

\begin{lem}[Lema Newman]
Yen $\red$ nduweni normalisasi kuwat lan sipat Church-Rosser ringkih,
relasi iki konfluen.
\end{lem}

\begin{proof}
  Amarga $\red$ nduweni normalisasi kuwat, saben urutan reduksi
  ngasilake wujud normal. Wujud normal tunggal yen lan mung yen $\red$
  konfluen. Mula upamane sawijining terma $N$ nduweni rong wujud normal
  sing beda, yaiku $N'$ lan $N''$, minangka asil saka urutan reduksi
  \begin{align*}
  N & = N_0' \redone N_1' \redone N_2' \redone \dots \redone N_k' = N' \text{ lan}\\
  N & = N_0'' \redone N_1'' \redone N_2'' \redone \dots \redone N_l'' = N''
  \end{align*}
  (Urutan reduksi kudu nduweni dawa $>0$, amarga yen ora, $N$ wis
  awujud normal.)

  Yen $N_1' = N_1''$, terma iki nduweni rong wujud normal sing beda
  ($N'$ lan $N''$). Yen $N_1' \neq N_1''$, miturut sipat Church-Rosser
  ringkih ana asil reduksi bebarengan. Kanthi normalisasi kuwat, asil
  kasebut nduweni wujud normal $N'''$, mula $N_1' \red N'''$ lan
  $N_1'' \red N'''$. Amarga $N' \neq N''$, salah siji saka
  $N''' \neq N'$ utawa $N''' \neq N''$ bener. Mula salah siji saka
  $N_1'$ utawa $N_1''$ nduweni rong wujud normal sing beda. Kita wis
  nuduhake yen nalika $N$ nduweni rong wujud normal sing beda, ana
  $N_1$ kanthi $N \redone N_1$ lan $N_1$ nduweni rong wujud normal
  sing beda. Kita bisa nerusake lan ngasilake $N \redone N_1
  \redone N_2 \redone \dots$, nanging iki mokal amarga $\redone$
  nduweni normalisasi kuwat.
\end{proof}


\end{document}
